% El ciclo de mejora — probar, hallar la falla, cambiar UNA cosa, volver a probar.
%
% Diagrama propio en TikZ (sin librerías extra: el círculo y las flechas curvas
% se dibujan a mano). Se tiñe con el color de la línea (milcGradeDark = naranja
% de Robótica). Muestra el ciclo de depuración/iteración como un lazo cerrado,
% con el principio de oro en el centro: cambiar una sola cosa por vuelta, para
% saber qué fue lo que la arregló.
%
% Se incrusta vía \guiaDiagrama desde el bloque `recursos:` del YAML.
\begin{tikzpicture}[scale=1,font=\sffamily,>=stealth]

  % Radio y centro del ciclo
  \def\R{2.4}
  \coordinate (C) at (0,0);

  % ── Cuatro estaciones del ciclo (explícitas para fijar el color de texto) ──
  \newcommand{\estacion}[5]{% ángulo · relleno · color-texto · título · detalle
    \node[fill=#2,rounded corners=3pt,inner sep=3pt,align=center,text=#3,
          font=\scriptsize\bfseries,text width=2.5cm]
      (S#1) at ({#1}:\R) {#4\\[0.4mm]\mdseries\itshape\footnotesize #5};
  }
  \estacion{90}{milcGradeDark}{white}{1. probar}{correr los 5 escenarios}
  \estacion{0}{milcMostaza}{milcNegro}{2. observar}{¿dónde falló? esperado vs real}
  \estacion{270}{milcGradeDark}{white}{3. cambiar UNA cosa}{un umbral o una regla}
  \estacion{180}{milcMostaza}{milcNegro}{4. volver a probar}{¿mejoró? anota el cambio}

  % ── Flechas curvas en el sentido del ciclo (horario) ────────────────
  \draw[->,milcGradeDark,line width=1pt] (S90) to[bend left=28] (S0);
  \draw[->,milcGradeDark,line width=1pt] (S0) to[bend left=28] (S270);
  \draw[->,milcGradeDark,line width=1pt] (S270) to[bend left=28] (S180);
  \draw[->,milcGradeDark,line width=1pt] (S180) to[bend left=28] (S90);

  % ── Centro: el principio de oro ──────────────────────────────────────
  \node[align=center,font=\scriptsize\bfseries,milcNegro,text width=2.7cm] at (C)
    {una sola cosa\\por vuelta\\[0.6mm]\mdseries\itshape\footnotesize (si cambias dos, no sabrás cuál la arregló)};

  % (La síntesis «mejorar es girar este ciclo…» va en el pie de figura vía
  %  \guiaDiagrama, para no repetirla dentro del propio diagrama.)

\end{tikzpicture}
